\section{截断、枚举与主定理的组合}\label{sec:conversion}
本节说明如何把集合内整除链结论转成引言中的正整数枚举形式。这里先处理数学对象，形式化声明的对应关系留到最后一节。

\subsection{端点与极限值域}
严格截断 $n<x$ 和非严格截断 $n\le x$ 只在 $x$ 是整数且该整数属于所计集合时有差别。对充分大的 $x$，有
\[
 0\le H_A^{\le}(x)-H_A(x)\le1,\qquad
 0\le W_A^{\le}(x)-W_A(x)\le\frac1{2\log2}.
\]
任意严格递增链的两种截断计数也至多相差 $1$。分别除以 $\log x$ 或 $\log\log x$ 后，这些误差均趋于零。因此倒数和的下极限、加权和的上极限及链计数的上极限都不受端点改变的影响。

加权密度由~\eqref{eq:bridge}~中的全体整数上界控制，是有限实数；链计数的归一化上极限则可能为 $+\infty$。例如 $1,2,4,8,\ldots$ 的计数为 $\log x/\log2+O(1)$。故主定理的计数上极限应放在 $[0,+\infty]$ 中解释。对最终非负且有界的加权比值，嵌入扩展非负实数与取上极限相容；计数一侧不需要额外假设有界。

\subsection{正整数子序列}
\begin{proof}[定理~\ref{thm:main}~的证明（使用已有整除链定理）]
令 $A=\{a_i:i\ge1\}$。严格递增性保证枚举没有重复，所以有限截断下的集合求和等于指标求和。由命题~\ref{prop:bridge}，$\Delta(A)>0$。应用定理~\ref{thm:upstream}，得到 $A$ 中严格递增的链 $(b_j)_{j\ge0}$ 及~\eqref{eq:upstream-bound}。

对每个 $i\ge1$，由 $b_{i-1}\in A$ 及枚举的单射性，存在唯一正整数 $k_i$ 使 $a_{k_i}=b_{i-1}$。如果 $k_i\ge k_{i+1}$，则 $a_{k_i}\ge a_{k_{i+1}}$，与 $b_{i-1}<b_i$ 矛盾。因此 $k_i<k_{i+1}$，且
\[
 a_{k_i}=b_{i-1}\mid b_i=a_{k_{i+1}}.
\]
指标平移 $i\mapsto i-1$ 给出相同的截断计数。再将~\eqref{eq:upstream-bound}~的两侧改为严格截断，即得~\eqref{eq:target}。
\end{proof}

在这条推导中，密度桥梁只负责提供调用上游定理所需的正性。最终计数下界是原集合的完整加权上密度 $\Delta(A)$，而不是桥梁证明中选取的辅助常数 $c$ 或 $c/2$。
